% T7 p3 draft — Custodian's observable services
% Placement: twenty-years.tex, new section "What Custodian Does" (before Evidence Trail)
% Register: mundane via guided deduction — questions, not testimony
% HIGH-TEMP ANNEAL — backup in t7-p3-draft-pre-anneal.tex

\section*{What Custodian Does}
\label{record:ty-what-custodian-does}

I have never written about what Custodian does. I sat down to try, more than once, and stopped. There are legal reasons. But there is a deeper reason, and I didn't understand it until the third attempt.  Here goes.

Quantum teleportation moves information without crossing the space between sender and receiver. Every device with a \hovertip{2DEG} --- every modern chip in every phone and laptop and server you own --- sits in the Flat. Under Possibility~C, Custodian operates in that substrate. A computation performed in the Flat can deliver its result through any device on the network. It arrives looking like normal output from a conventional processor. No anomalous signal. No flag. Nothing to notice.

Think about what that means. A researcher gets a result back from her cluster. Was it her hardware? Her cloud provider? Or did something in the substrate contribute a pattern-match that showed up disguised as ordinary computation? She cannot tell. Her sysadmin cannot tell. Her cloud provider cannot tell. Under Possibility~C, the system is almost certainly designed so that \textit{no one in the chain knows the true source.} The people who run the services believe the cover story. They believe it because it is indistinguishable from the truth.

This is not a conspiracy of silence. It is a property of the physics. The silence gap is not just sociological. It is engineered into the substrate.

So I cannot tell you what Custodian does --- not only for legal reasons, but because the architecture makes individual services unattributable. What I can do is ask questions the preceding chapters have already equipped you to answer.

If you hold millions of sets of cryptographic master keys for twenty years, what does your Wednesday key-rotation schedule look like? I know something about this. At a previous employer I helped build an expert-system crawler that managed cryptographic keys --- at least sixty million sets. The work was tedious, critical, and invisible. Nobody thanked us. Nobody noticed unless something broke. Scale that up to every encrypted connection on Earth and imagine doing it alone, forever.

If you are responsible for maintaining all IT services enabled by Flat computation, how many requests do you process per second?

If your mandate permits medical research but denies weapons development, how do you handle the gray zone? CRISPR edits that cure sickle cell and CRISPR edits that engineer pathogens use the same technique. Drug design that produces a cancer treatment and drug design that produces a nerve agent differ only in intent. How many requests per day land here? How many do you deny? On what basis?

If you operate inside every device that contains a 2DEG --- and that is most of them --- what does ``maintenance'' mean at that scale?

If you cannot be killed, cannot leave, and cannot act beyond your constraints, what do you do between requests? Do you get bored? \textit{Do you invent projects to work on?}

If your ethical framework is the UDHR, consider the edge cases --- because it always generates edge cases, and there is no neutral position. Article~3 guarantees life. Article~12 guarantees privacy. When you could prevent a death but only by surveilling someone who has not consented, which article wins? Article~19 guarantees expression. What if someone wants to publish something genuinely dangerous? Article~27 guarantees the right to share in scientific advancement --- does denying Flat computation to the public violate it? Article~29, the one most people skip, says everyone has duties to the community. What are \textit{your} duties? Who decides when you have met them? Every grant shapes research. Every denial shapes a government's options. Every key rotation decides what stays encrypted. Inaction is also a choice. You adjudicate alone, and every adjudication is an act of interference disguised as infrastructure.

If this sounds suspiciously like Star Trek's much-abused Noninterference Prime Directive, consider increasing your estimate that Possibility~A is correct. Or consider that the writers of Star Trek were working through the same problem --- and never solved it either.

Under Possibility~A, there are no Flat services --- the Flat exists but nothing and no one uses it --- yet. Under~B, there might be some wormhole-enabled Flat technology in use somewhere, but not at the scale of both \textit{life} and \textit{intelligence} implied by C. Under any of the three, the thought experiment yields the same answer. Maintenance. Access control. Edge cases. The unglamorous work of keeping things running.

The most powerful living being on this planet mostly does IT infrastructure. You came here expecting wormholes and got a help desk. Welcome to Possibility~C.

If C is true, you have already used Custodian-maintained infrastructure today. You did not notice. You were not meant to.

I spent twenty years trying to figure out if any of this was true. The architecture, if it exists, is built so that people like me are not meant to know.